\documentclass[11pt, letterpaper]{article}
\usepackage{mobi-lectura}

\title{Cadenas de Markov de Tiempo Continuo: \\ Del Enunciado a la Matriz Generadora}
\author{Alejandra Tabares Pozos}
\date{}

\begin{document}

\maketitle
\tableofcontents
\newpage

% ══════════════════════════════════════════════════
\section{Conceptos esenciales}
% ══════════════════════════════════════════════════

Una CTMC $\{X(t), t\ge0\}$ queda completamente descrita por su \textbf{matriz generadora} $\mat{Q}=[q_{ij}]$, donde $q_{ij}\ge0$ ($i\neq j$) es la tasa de transici\'on de $i$ a $j$, y $q_{ii}=-\sum_{j\neq i}q_{ij}$ (filas suman cero). El tiempo de permanencia en el estado $i$ es $\text{Exp}(|q_{ii}|)$.

\begin{metodobox}[Metodolog\'ia en 5 pasos]
\begin{enumerate}
    \item \textbf{Variables aleatorias:} \textquestiondown{}qu\'e cantidades cambian en el tiempo?
    \item \textbf{Espacio de estados:} \textquestiondown{}qu\'e valores puede tomar cada variable?
    \item \textbf{Estado conjunto:} si hay varias variables, el estado es una tupla; el espacio es el producto cartesiano.
    \item \textbf{Transiciones:} para cada evento, definir el cambio de estado, su tasa y cu\'ando es posible.
    \item \textbf{Matriz $\mat{Q}$:} ensamblar las entradas, de forma \emph{expl\'icita} (num\'erica) o \emph{impl\'icita} (funciones de los \'indices).
\end{enumerate}
\end{metodobox}

% ══════════════════════════════════════════════════
\section{Escenarios con una variable (procesos nacimiento--muerte)}
% ══════════════════════════════════════════════════

En estos sistemas, el estado es un entero $n$ y las transiciones solo van a $n\pm1$.

% ──────────────────────────────────────────────
\subsection{Escenario A: Cola M/M/1}

\begin{ejemplobox}[Enunciado]
Un cajero atiende clientes que llegan con tasa $\lambda=3$/min y son atendidos con tasa $\mu=5$/min. La fila no tiene l\'imite.
\end{ejemplobox}

\textbf{Variable:} $N(t)=$ clientes en el sistema. \textbf{Espacio:} $\{0,1,2,\ldots\}$.

\begin{center}
\begin{tikzpicture}[state/.style={circle, draw=uniblue, fill=uniblue!10, minimum size=9mm, font=\footnotesize\bfseries}, >=Stealth]
\node[state](0)at(0,0){0};\node[state](1)at(2.2,0){1};\node[state](2)at(4.4,0){2};\node[state](3)at(6.6,0){3};\node[font=\large]at(8.2,0){$\cdots$};
\foreach \from/\to in {0/1,1/2,2/3}{
    \draw[->,thick,uniblue](\from)to[bend left=25]node[above,font=\scriptsize]{$\lambda$}(\to);
    \draw[->,thick,red!70!black](\to)to[bend left=25]node[below,font=\scriptsize]{$\mu$}(\from);
}
\end{tikzpicture}
\end{center}

\textbf{Forma impl\'icita:} \quad $q_{n,n+1}=\lambda$, \quad $q_{n,n-1}=\mu\,\ind{n\ge1}$, \quad $q_{nn}=-\lambda-\mu\,\ind{n\ge1}$.

\textbf{Matriz expl\'icita} (primeras filas):
$\mat{Q}=\begin{psmallmatrix}-3&3&0&\cdots\\5&-8&3&\cdots\\0&5&-8&\cdots\\\vdots&&&\ddots\end{psmallmatrix}$

% ──────────────────────────────────────────────
\subsection{Escenario B: Cola M/M/1/K (capacidad finita)}

\begin{ejemplobox}[Enunciado]
Cl\'inica dental: 1 dentista, capacidad total $K=4$. Llegadas $\lambda=2$/h, servicio $\mu=3$/h. Si el sistema est\'a lleno, el paciente se va.
\end{ejemplobox}

La \'unica diferencia con la M/M/1 es que las llegadas se bloquean cuando $n=K$:

\[
q_{n,n+1}=\lambda\,\ind{n<K}, \qquad q_{n,n-1}=\mu\,\ind{n\ge1}
\]

\[
\mat{Q}=\begin{pmatrix}-2&2&0&0&0\\3&-5&2&0&0\\0&3&-5&2&0\\0&0&3&-5&2\\0&0&0&3&-3\end{pmatrix}
\]

Observe: la primera fila no tiene $\mu$ (sistema vac\'io) y la \'ultima no tiene $\lambda$ (sistema lleno).

% ──────────────────────────────────────────────
\subsection{Escenario C: Cola M/M/c (m\'ultiples servidores)}

\begin{ejemplobox}[Enunciado]
Centro de llamadas: $c=3$ operadores, llegadas $\lambda=10$/min, servicio por operador $\mu=4$/min, cola infinita.
\end{ejemplobox}

La tasa de servicio \textbf{depende del estado}: si hay $n$ llamadas, hay $\min(n,c)$ operadores activos.

\[
q_{n,n+1}=\lambda, \qquad q_{n,n-1}=\min(n,c)\cdot\mu, \qquad q_{nn}=-\lambda-\min(n,c)\cdot\mu
\]

\[
\mat{Q}=\begin{pmatrix}-10&10&0&0&0&\cdots\\4&-14&10&0&0&\cdots\\0&8&-18&10&0&\cdots\\0&0&12&-22&10&\cdots\\0&0&0&12&-22&\cdots\\\vdots&&&&&\ddots\end{pmatrix}
\]

La tasa de muerte crece ($\mu, 2\mu, 3\mu$) hasta saturarse en $c\mu=12$ a partir de $n=c$.

\begin{notabox}[Patr\'on general: nacimiento--muerte]
Los tres escenarios comparten la estructura $q_{n,n+1}=\lambda_n$, $q_{n,n-1}=\mu_n$. Lo que los distingue es la forma funcional:
\begin{center}\small
\begin{tabular}{lcc}\toprule
\textbf{Modelo} & $\lambda_n$ & $\mu_n$ \\\midrule
M/M/1 & $\lambda$ & $\mu$ \\
M/M/1/K & $\lambda\,\ind{n<K}$ & $\mu\,\ind{n\ge1}$ \\
M/M/c & $\lambda$ & $\min(n,c)\,\mu$ \\\bottomrule
\end{tabular}
\end{center}
\end{notabox}

% ══════════════════════════════════════════════════
\section{Escenarios con dos variables}
% ══════════════════════════════════════════════════

Cuando una sola variable no captura toda la informaci\'on necesaria para determinar las tasas de transici\'on, necesitamos m\'ultiples variables. El estado es una tupla y el espacio es un producto cartesiano.

% ──────────────────────────────────────────────
\subsection{Escenario D: Reparaci\'on de m\'aquinas}

\begin{ejemplobox}[Enunciado]
$M=2$ m\'aquinas id\'enticas, $R=1$ t\'ecnico. Tasa de falla $\alpha=0.5$/h, tasa de reparaci\'on $\beta=2$/h. Si ambas fallan, el t\'ecnico repara una a la vez.
\end{ejemplobox}

Podemos modelar con dos variables $X_i(t)\in\{0,1\}$ (estado de cada m\'aquina), dando un espacio $\{(0,0),(0,1),(1,0),(1,1)\}$ con 4 estados. Pero como las m\'aquinas son \textbf{id\'enticas}, reducimos a una sola variable:
\[
F(t) = \text{m\'aquinas da\~nadas} \in \{0, 1, 2\}
\]

Tasas: falla con tasa $(M-f)\alpha$ (m\'aquinas operativas) y reparaci\'on con tasa $\min(f,R)\beta$ (t\'ecnicos activos).

\begin{center}
\begin{tikzpicture}[state/.style={circle, draw=uniblue, fill=uniblue!10, minimum size=10mm, font=\footnotesize\bfseries}, >=Stealth]
\node[state](0)at(0,0){0};\node[state](1)at(3.5,0){1};\node[state](2)at(7,0){2};
\draw[->,thick,uniblue](0)to[bend left=25]node[above,font=\scriptsize]{$2\alpha$}(1);
\draw[->,thick,uniblue](1)to[bend left=25]node[above,font=\scriptsize]{$\alpha$}(2);
\draw[->,thick,red!70!black](1)to[bend left=25]node[below,font=\scriptsize]{$\beta$}(0);
\draw[->,thick,red!70!black](2)to[bend left=25]node[below,font=\scriptsize]{$\beta$}(1);
\end{tikzpicture}
\end{center}

\textbf{Forma impl\'icita general} ($M$ m\'aquinas, $R$ t\'ecnicos):
\[
\boxed{q_{f,f+1}=(M-f)\alpha, \qquad q_{f,f-1}=\min(f,R)\,\beta}
\]

Con $M=2, R=1, \alpha=0.5, \beta=2$: \quad
$\mat{Q}=\begin{pmatrix}-1&1&0\\2&-2.5&0.5\\0&2&-2\end{pmatrix}$

% ──────────────────────────────────────────────
\subsection{Escenario E: Red de colas con bloqueo (dos variables irreducibles)}

\begin{ejemplobox}[Enunciado]
Taller con dos estaciones en serie (A: diagn\'ostico, B: reparaci\'on). Capacidad m\'axima 2 en cada una. Llegadas $\lambda=2$/h, servicio $\mu_A=4$/h, $\mu_B=3$/h. Si B est\'a llena, A queda \textbf{bloqueada}.
\end{ejemplobox}

Aqu\'i \textbf{no podemos reducir} a una variable: las tasas dependen de c\'omo se distribuyen las bicicletas entre A y B, no solo del total.

\textbf{Variables:} $N_A(t)\in\{0,1,2\}$, $N_B(t)\in\{0,1,2\}$. \textbf{Estado:} $(a,b)$, con $|\mathcal{S}|=9$.

\textbf{Tabla de eventos:}

\begin{center}\small
\renewcommand{\arraystretch}{1.2}
\begin{tabular}{lccc}\toprule
\textbf{Evento} & \textbf{Cambio $\vect{\delta}$} & \textbf{Tasa} & \textbf{Condici\'on} \\\midrule
Llegada & $(+1,\,0)$ & $\lambda$ & $a<2$ \\
Servicio A$\to$B & $(-1,\,+1)$ & $\mu_A$ & $a\ge1$ \textbf{y} $b<2$ \\
Salida de B & $(0,\,-1)$ & $\mu_B$ & $b\ge1$ \\\bottomrule
\end{tabular}
\end{center}

El bloqueo aparece en la segunda fila: A solo puede despachar si B tiene espacio.

\textbf{Forma impl\'icita:}
\[
\boxed{\begin{aligned}
q_{(a,b),(a+1,b)} &= \lambda\,\ind{a<2} \\
q_{(a,b),(a-1,b+1)} &= \mu_A\,\ind{a\ge1}\,\ind{b<2} \\
q_{(a,b),(a,b-1)} &= \mu_B\,\ind{b\ge1}
\end{aligned}}
\]

\textbf{Diagrama de transiciones:}

\begin{figure}[H]
\centering
\begin{tikzpicture}[
    state/.style={circle, draw=uniblue, fill=uniblue!10, minimum size=8mm, font=\tiny\bfseries, inner sep=0pt},
    >=Stealth, scale=0.88, transform shape]
\foreach\a in{0,1,2}{\foreach\b in{0,1,2}{\node[state](\a\b)at(3.2*\a,-2*\b){$\a,\b$};}}
\foreach\a/\an in{0/1,1/2}{\foreach\b in{0,1,2}{\draw[->,thick,uniblue](\a\b)--node[above,font=\tiny]{$\lambda$}(\an\b);}}
\foreach\a/\ap in{1/0,2/1}{\foreach\b/\bn in{0/1,1/2}{\draw[->,thick,darkgreen!80](\a\b)--node[above,sloped,font=\tiny]{$\mu_A$}(\ap\bn);}}
\foreach\a in{0,1,2}{\foreach\b/\bp in{1/0,2/1}{\draw[->,thick,red!70!black](\a\b)--node[right,font=\tiny]{$\mu_B$}(\a\bp);}}
\node[above=4mm,font=\small\bfseries,uniblue]at(00.north){$a{=}0$};
\node[above=4mm,font=\small\bfseries,uniblue]at(10.north){$a{=}1$};
\node[above=4mm,font=\small\bfseries,uniblue]at(20.north){$a{=}2$};
\node[left=4mm,font=\small\bfseries,red!70!black]at(00.west){$b{=}0$};
\node[left=4mm,font=\small\bfseries,red!70!black]at(01.west){$b{=}1$};
\node[left=4mm,font=\small\bfseries,red!70!black]at(02.west){$b{=}2$};
\end{tikzpicture}
\caption{Red de colas: flechas $\textcolor{uniblue}{\rightarrow}$ llegadas, $\textcolor{darkgreen!80}{\searrow}$ servicio A, $\textcolor{red!70!black}{\uparrow}$ servicio B. No hay flechas $\mu_A$ desde la fila $b=2$ (bloqueo).}
\end{figure}

\textbf{Matriz expl\'icita} (orden lexicogr\'afico):
\[
\mat{Q} = \left(\!\begin{smallmatrix}
-2 & 0 & 0 & 2 & 0 & 0 & 0 & 0 & 0 \\
3 & -5 & 0 & 0 & 2 & 0 & 0 & 0 & 0 \\
0 & 3 & -3 & 0 & 0 & 2 & 0 & 0 & 0 \\
4 & 0 & 0 & -6 & 0 & 0 & 2 & 0 & 0 \\
0 & 4 & 0 & 3 & -9 & 0 & 0 & 2 & 0 \\
0 & 0 & 0 & 0 & 3 & -5 & 0 & 0 & 2 \\
0 & 0 & 0 & 4 & 0 & 0 & -4 & 0 & 0 \\
0 & 0 & 0 & 0 & 4 & 0 & 3 & -7 & 0 \\
0 & 0 & 0 & 0 & 0 & 0 & 0 & 3 & -3
\end{smallmatrix}\right)
\]

% ══════════════════════════════════════════════════
\section{Modelo expl\'icito vs.\ impl\'icito}
% ══════════════════════════════════════════════════

\begin{center}\small
\renewcommand{\arraystretch}{1.3}
\begin{tabularx}{\textwidth}{>{\bfseries}l X X}\toprule
& \textbf{Expl\'icito} & \textbf{Impl\'icito} \\\midrule
Qu\'e es & Matriz $\mat{Q}$ con valores num\'ericos & Tasas como funciones de los \'indices \\
Ventaja & Visual, verificable a mano & Compacto, generalizable, programable \\
Usar cuando & $|\mathcal{S}|\le 10$ o para verificar & $|\mathcal{S}|$ grande o parametrizado \\\bottomrule
\end{tabularx}
\end{center}

El patr\'on general para la forma impl\'icita es: para cada tipo de evento $k$, definir un \textbf{vector de cambio} $\vect{\delta}_k$, una \textbf{tasa} $r_k(\vect{x})$, y una \textbf{condici\'on} $C_k(\vect{x})$. Entonces:
\[
q_{\vect{x},\,\vect{x}+\vect{\delta}_k} = r_k(\vect{x})\cdot\ind{C_k(\vect{x})}, \qquad q_{\vect{x},\vect{x}} = -\sum_k r_k(\vect{x})\cdot\ind{C_k(\vect{x})}
\]

\begin{notabox}[Resumen de todos los escenarios]
\begin{center}\footnotesize
\renewcommand{\arraystretch}{1.2}
\begin{tabular}{lccccc}\toprule
& \textbf{A: M/M/1} & \textbf{B: M/M/1/K} & \textbf{C: M/M/c} & \textbf{D: Reparac.} & \textbf{E: Red} \\\midrule
Variables & 1 & 1 & 1 & 1 (reducida) & 2 \\
$|\mathcal{S}|$ & $\infty$ & $K{+}1$ & $\infty$ & $M{+}1$ & $(c_A{+}1)(c_B{+}1)$ \\
Tasas & ctes. & ctes.+borde & dep.\ estado & dep.\ estado & dep.\ estado \\
Acople & no & no & no & no & s\'i (bloqueo) \\\bottomrule
\end{tabular}
\end{center}
\end{notabox}

% ══════════════════════════════════════════════════
\section{Conexión con el análisis transitorio y estacionario}
% ══════════════════════════════════════════════════

Una vez construida la matriz $\mat{Q}$, el siguiente tema estudia qué hace el sistema a lo largo del tiempo. Existen dos preguntas naturales:

\begin{enumerate}
    \item \textbf{Comportamiento transitorio:} ¿cuál es la distribución de $X(t)$ en el instante $t$? La respuesta involucra la \emph{exponencial matricial} $e^{\mat{Q}t}$: si $\vect{\pi}(0)$ es la distribución inicial, entonces
    \[
    \vect{\pi}(t) = \vect{\pi}(0)\, e^{\mat{Q}t}.
    \]

    \item \textbf{Distribución estacionaria $\vect{\pi}^*$:} si la cadena es ergódica, cuando $t\to\infty$ la distribución converge a un límite independiente del estado inicial. Este límite satisface el sistema lineal:
    \[
    \boxed{\vect{\pi}^* \mat{Q} = \vect{0}, \qquad \textstyle\sum_i \pi^*_i = 1.}
    \]
\end{enumerate}

\textbf{Ejemplo: M/M/1/K del Escenario B} ($\lambda=2$, $\mu=3$, $K=4$).

El sistema $\vect{\pi}^*\mat{Q}=\vect{0}$ produce cinco ecuaciones de balance:
\begin{align*}
-2\pi_0 + 3\pi_1 &= 0 & 2\pi_0 - 5\pi_1 + 3\pi_2 &= 0 \\
2\pi_1 - 5\pi_2 + 3\pi_3 &= 0 & 2\pi_2 - 5\pi_3 + 3\pi_4 &= 0 & 2\pi_3 - 3\pi_4 &= 0
\end{align*}
junto con $\sum_{i=0}^4 \pi_i = 1$. Por sustitución progresiva, $\pi_i = \pi_0\,(2/3)^i$. De la normalización:
\[
\pi_0\sum_{i=0}^4 (2/3)^i = 1 \implies \pi_0 = \frac{1}{1+\tfrac{2}{3}+\tfrac{4}{9}+\tfrac{8}{27}+\tfrac{16}{81}} \approx 0{,}417.
\]
La distribución estacionaria es $(\pi_0,\pi_1,\pi_2,\pi_3,\pi_4) \approx (0{,}417,\;0{,}278,\;0{,}185,\;0{,}123,\;0{,}082)$. Interpretación: el sistema está vacío el 41.7\% del tiempo; la tasa efectiva de llegadas es $\lambda_{\text{ef}} = \lambda(1-\pi_4) \approx 2(0{,}918) = 1{,}836$ clientes/h.

\begin{alertabox}[Para el próximo tema]
La construcción correcta de $\mat{Q}$ es el prerrequisito indispensable. Un error en las tasas o en las condiciones de borde se propaga directamente a $e^{\mat{Q}t}$ y a $\vect{\pi}^*\mat{Q}=\vect{0}$. Verifique siempre que cada fila de $\mat{Q}$ sume cero antes de continuar.
\end{alertabox}

\vfill

\section*{Referencias}
\begin{enumerate}[label={[\arabic*]}]
    \item Ross, S.\,M. (2014). \textit{Introduction to Probability Models}, 11th ed. Academic Press. Cap.~6.
    \item Stewart, W.\,J. (2009). \textit{Probability, Markov Chains, Queues, and Simulation}. Princeton University Press. Cap.~8.
    \item Gross, D., Shortle, J., Thompson, J., \& Harris, C. (2008). \textit{Fundamentals of Queueing Theory}, 4th ed. Wiley. Cap.~3--4.
    \item Karlin, S. \& Taylor, H.\,M. (1975). \textit{A First Course in Stochastic Processes}, 2nd ed. Academic Press. Cap.~4.
    \item \c{C}inlar, E. (2013). \textit{Introduction to Stochastic Processes}. Dover. Cap.~8.
\end{enumerate}

\end{document}
